% Purpose:
% Frame the paper as an ED-track contribution, not a systems or solver paper.
%
% Must do:
% - Introduce ARC as a benchmark increasingly studied through test-time adaptation and compute.
% - Explain why systematic support-to-query evaluation is hard without a practical executable substrate.
% - Introduce JaxARC as the enabler of that evaluation.
% - Preview the central empirical finding: support-side gains do not become hidden-query success.
% - End with a crisp contributions list.
%
% Keep out:
% - Long API or implementation detail.
% - Overstated speed bragging.
% - Any implication that the baselines solve hidden queries.
%
% Primary sources:
% - planning/neurips_ed_paper_writing_guide.md
% - planning/neurips_ed_final_implementation_plan.md
% - results/neurips_ed_final_program/STATUS_SUMMARY.md

The Abstraction and Reasoning Corpus (ARC) tests an AI system's ability to infer rules from a few examples \citep{cholletMeasureIntelligence2019}. While most approaches treat ARC as a static problem solved via program synthesis \citep{alfordNeuralGuidedBidirectionalProgram2022} or large language models \citep{buttCodeItSelfimprovingLanguage2024,batorskiNSANeurosymbolicARC2025}, cognitive science suggests that abstract reasoning depends heavily on sensorimotor experience \citep{thelenGroundedWorldDevelopmental2000,zadorCatalyzingNextgenerationArtificial2023}. Biological intelligence relies on continuous interaction with the environment to acquire new skills \citep{agrawalComputationalSensorimotorLearning2018}. By reducing ARC to static input-output pairs, we strip away the interactive feedback loops that ordinarily support learning.

Formulating ARC as a reinforcement learning (RL) problem restores this interaction. An agent starts with an input grid, applies sequential edits, observes the resulting state, and eventually submits an answer. ARCLE \citep{leeARCLEAbstractionReasonin2024} introduced this perspective by building a Python-based ARC environment and evaluating baseline PPO agents. In this framing, ARC becomes a testbed for long-horizon credit assignment and generalization across diverse task families.

But scaling these experiments is difficult. Modern RL relies on hardware accelerators and highly parallelized environments. When the simulation logic runs on the CPU, as it does in existing Python environments, the environment becomes the primary bottleneck. Standard sweeps over hyperparameters, tasks, and network architectures take weeks. We have seen similar bottlenecks in grid-world navigation until JAX-based reimplementations made large-scale research feasible again \citep{pignatelliNAVIXScalingMiniGrid2025}. ARC requires the same transition to make task-specific adaptation studies practical.

To address this, we introduce JaxARC, a JAX-native environment that runs ARC entirely on accelerators. It provides a functional API, fixed-shape buffers, and native support for major datasets like ARC-AGI \citep{cholletARCAGI2NewChallenge2025}, ConceptARC \citep{moskvichevConceptARCBenchmarkEvaluating2023}, and MiniARC \citep{kimPlaygroundsAbstractionReasoning2022}. Following the design of recent JAX environments \citep{koyamadaPgxHardwareAcceleratedParallel2023,bonnetJumanjiDiverseSuite2023}, we prioritize throughput to make large-scale experiments viable. On an H100 GPU, JaxARC executes over 100 million steps per second. We profile this throughput and measure the computational overhead of various observation and action wrappers to map the environment's scaling properties.

Fast simulation enables more rigorous evaluation protocols. Because every ARC task provides "support" examples for learning and "query" examples for testing, an agent can easily fit the support grids without solving the underlying rule. We built a compute-budgeted benchmark on top of JaxARC that strictly separates support-side adaptation from hidden-query transfer. When we evaluate canonical feedforward and recurrent PPO baselines under this protocol, the results show a clear dissociation. Agents reliably learn to solve the support grids as their compute budget increases, but their success rate on hidden queries remains exactly zero. In our ConceptARC runs, the majority of tasks reach a support similarity of 0.8 or higher while still failing the queries completely.

We make the following contributions:

\begin{itemize}
	\item We release JaxARC, an open-source, JAX-native RL environment for ARC that operates at over 100 million steps per second.
	\item We define a compute-budgeted evaluation protocol that isolates task-specific adaptation from hidden-query generalization.
	\item We provide comprehensive throughput profiling and baseline evaluations showing that current PPO methods successfully fit ARC support demonstrations but fail to transfer that behavior to hidden queries.
\end{itemize}

JaxARC makes ARC much easier to study as an RL domain. By doing so, it exposes exactly where current adaptation methods stop short and establishes a clear, measurable gap for future research.